\documentclass[a4paper, 14pt]{article}

\usepackage{cmap}
\usepackage[T2A]{fontenc}
\usepackage[utf8]{inputenc}
\usepackage[russian]{babel}
\usepackage[left=3cm,right=1.5cm,top=2cm,bottom=2cm]{geometry}
\usepackage{indentfirst}
\linespread{1.5}
\setlength{\parindent}{1.25cm}
\pagestyle{empty}

\newcommand{\fieldline}[1]{\noindent\underline{\parbox[t]{\dimexpr\linewidth-2pt\relax}{#1}}\par}

\begin{document}

\begin{center}
    \small
    Приложение~4\\
    к Положению о выпускной\\
    квалификационной работе студентов МФТИ\\[1em]

    \normalsize
    Рекомендуемая форма отзыва научного руководителя на ВКР\\[0.5em]
    \textit{Оформляется в печатном виде}\\[1.5em]

    \Large\textbf{ОТЗЫВ НАУЧНОГО РУКОВОДИТЕЛЯ}\\[0.5em]
    \large\textbf{на выпускную квалификационную работу}
\end{center}

\vspace{1em}

\noindent\fieldline{\parbox[t]{\linewidth}{%
    \textbf{Абубакиров Владимир Андреевич}\\
    \textit{«Применение многопоточного кольцевого буфера для организации высокоскоростного обмена диагностической информацией в среде виртуализации»}%
}}

\vspace{1.5em}

\noindent Направление подготовки / специальность\\
\fieldline{01.03.02 «Прикладная математика и информатика»}

\vspace{1em}

\noindent Направленность (профиль) подготовки\\
\fieldline{«Математика и компьютерные науки»}

\vspace{1.5em}

\textbf{Актуальность исследования.}
Развитие облачных платформ и повсеместное использование виртуальных машин
повышают требования к инструментам диагностики и отладки гостевых операционных
систем. Традиционные каналы передачи логов связаны с существенными накладными
расходами и не всегда обеспечивают достаточную пропускную способность при
интенсивной генерации событий. Дополнительную сложность представляет сбор
диагностики на ранних этапах загрузки (UEFI) и из различных компонентов гостевой
системы --- приложений, сервисов и драйверов. Тема работы актуальна для задач
разработки и сопровождения виртуализированных сред.

\textbf{Основные результаты исследования.}
В работе разработан комплекс программных компонентов для высокоскоростной
передачи диагностической информации из гостевой системы в эмулятор Stratovirt:
(1)~lock-free MPSC-кольцевой буфер в общей памяти с алгоритмом резервирования,
копирования и фиксации записи без блокирующих примитивов синхронизации;
(2)~виртуальное устройство LogDev на стороне гипервизора с доступом через MMIO
и ACPI, обеспечивающее запись сообщений в файл на хосте;
(3)~UEFI-драйвер и адаптация DebugLib для логирования на этапах прошивки и после
вызова \texttt{ExitBootServices};
(4)~драйвер операционной системы Windows с интерфейсами доступа к буферу из
пользовательского и ядерного режима.
Предложенное решение обеспечивает единый хронологически упорядоченный поток
диагностических сообщений от гостевой ОС и эмулятора.

\textbf{Достоверность результатов.}
Алгоритм буфера описан с учётом модели памяти и механизмов целостности записей;
приведены схемы состояний и псевдокод ключевых операций. Архитектура компонентов
и их взаимодействие согласованы с реализацией в Stratovirt, EDK2 и Windows Driver
Framework. Выводы в заключении подкреплены сравнением с блокирующими и
неблокирующими подходами к организации кольцевых буферов.

\textbf{Теоретическая и практическая значимость.}
Теоретический вклад работы состоит в применении lock-free MPSC-кольцевого буфера
как канала связи «гость~---~хост» в виртуализированной среде с учётом особенностей
планирования виртуальных процессоров и конкурентной записи из нескольких потоков.
Практическая значимость определяется готовым прототипом, интегрируемым в Stratovirt
и применимым при отладке гостевых систем Windows и прошивки UEFI.

\textbf{Положительные стороны и недостатки исследования.}
К положительным сторонам следует отнести логичную структуру работы, последовательное
раскрытие архитектуры решения и продуманную сквозную интеграцию компонентов --- от
прошивки до прикладного программного обеспечения. Студент продемонстрировал владение
низкоуровневым системным программированием и корректное использование терминологии
виртуализации, UEFI и WDF. К недостаткам можно отнести отсутствие экспериментальной
оценки производительности и количественного сравнения с альтернативными каналами
передачи диагностики, а также необходимость доработки отдельных формулировок и
иллюстративного псевдокода.

\textbf{Характеристика работы студента над темой.}
Абубакиров~В.~А. самостоятельно изучил документацию Stratovirt, UEFI EDK2 и Windows
Driver Kit, проанализировал существующие подходы к передаче диагностики в среде
виртуализации. Объём проанализированного материала достаточен для задач бакалаврской
работы. Студент последовательно реализовал все заявленные программные компоненты,
умело работал с технической литературой и исходным кодом открытых проектов, уверенно
излагает архитектуру решения и отвечает на вопросы по lock-free алгоритмам и
интеграции драйверов.

\textbf{Оценка компетентности студента как будущего бакалавра.}
Работа показывает, что студент владеет навыками системного программирования,
понимает архитектуру виртуализации и драйверов, способен проектировать и реализовывать
нетривиальные механизмы конкурентного доступа к данным. Выпускник может претендовать на
квалификацию бакалавра по направлению подготовки.

\vspace{1em}

\noindent\textbf{Оценка выпускной работы}
(«отлично», «хорошо», «удовлетворительно», «неудовлетворительно»): \underline{\hspace{4cm}}.

Заключение о соответствии работы требованиям, установленным к ВКР студентов МФТИ,
и рекомендация о присвоении выпускнику соответствующей квалификации: выпускная
квалификационная работа соответствует установленным требованиям по объёму, структуре
и содержанию; рекомендую присвоить выпускнику квалификацию \textbf{бакалавра}.

\vspace{2em}

\noindent\textbf{Научный руководитель:}\\
кандидат физико-математических наук\\
Камкин Александр Сергеевич

\vspace{2em}

\noindent\underline{\hspace{6cm}}\\
\noindent\textit{(подпись научного руководителя)}

\vspace{2em}

\noindent\hfill «\underline{\hspace{1cm}}» \underline{\hspace{3cm}} 2026~г.

\end{document}
